\subsection{Exact scanning of predefined straight screenlines}\label{sec:scan}
Fix $D$ directions before evaluating attribution. For each direction, sort nodes by their coordinate projection and group equal projection values; no screenline splits such a group. A candidate position is a gap between adjacent groups. Forward and reverse union--find sweeps record whether each prefix and suffix induces a connected subgraph. This geometry and topology preprocessing is reusable across queries, demand paths and models.
For one attribution column, form $b_j=Bc_j$ (Eq.~\eqref{eq:Cdef}) and prefix sums $P_{j,k}=\sum_{i\leq k}(b_j)_{\pi(i)}$. If $s_j$ lies in the prefix and $t_j$ outside, its attribution is $-P_{j,k}$; for the reverse orientation it is $P_{j,k}$ because $\mathbf1^\top b_j=0$. Retain only endpoint-separating gaps with both sides connected. Scanning these values gives exact attributions and extrema over this prescribed finite screenline family. Preprocessing takes $O(D(|V|\log|V|+|E|\alpha(|V|)))$; after preprocessing, a column takes $O(|E|+D|V|)$ using sparse incidence multiplication. Duplicate regions may be removed without changing extrema.

On Sioux Falls, 36 directions at five-degree intervals give 1642 admissible query--position boundary values across six queries (duplicates across directions retained for all algorithms). On exactly these oriented regions, the prefix scan, a vectorized baseline that precomputes the signed crossing matrix $A_{\mathcal F}$ (one sparse row per oriented region, $+1$ leaving and $-1$ entering) and aggregates every boundary value with one sparse product $A_{\mathcal F}C$, and direct crossing-edge summation all agree to $2.2\times10^{-11}$ TSTT units once the reference columns are available. Appendix~\ref{supp:timing} reports the same-task timings of the three aggregations and the amortization threshold of the precomputed matrix.
